\documentclass{article}
\usepackage[utf8]{inputenc}
\usepackage{amsmath,amssymb}
\usepackage[most]{tcolorbox}
\usepackage{geometry}
\geometry{margin=2.5cm}

\newcommand{\step}[1]{\par\noindent\hangindent=3.5em\hangafter=1%
  \makebox[3.5em][l]{\textbf{#1.}}}
\newcommand{\steptag}[1]{\hfill{\small\textsf{[#1]}}}

\newtcolorbox{proofbox}[1]{
  colback=gray!5, colframe=gray!60,
  fonttitle=\bfseries, title={#1},
  breakable, enhanced,
  left=4pt, right=4pt, top=4pt, bottom=4pt,
  before skip=10pt, after skip=10pt
}

\begin{document}

\begin{proofbox}{Parameterized Maurer-Cartan transport: there exist C\_ch\textasciicircum{}0...}
\step{1} Parameterized Maurer-Cartan transport: there exist C\_ch\textasciicircum{}0 >= 1 and kappa\_ch\textasciicircum{}0 in (0, 1/2] such that, for every def-stage1-polar-witness-data tuple W with C\_ch >= C\_ch\textasciicircum{}0 and 0 < kappa\_ch <= kappa\_ch\textasciicircum{}0, for every finite-dimensional exact-unit epsilon\_r-C*-algebra, every delta > 0 with C\_ch*(epsilon\_r + delta) <= kappa\_ch, and every family g = (g\_U)\_\{U in calU\} of C\textasciicircum{}1 maps g\_U: B\textasciicircum{}\{icalH\}\_\{2delta\}(0) -> B\textasciicircum{}\{calH\}\_\{2delta\}(0) such that, for every U in calU and A\textasciicircum{}par in B\textasciicircum{}\{icalH\}\_\{2delta\}(0), g\_U(A\textasciicircum{}par) is the unique element of B\textasciicircum{}\{calH\}\_\{2delta\}(0) satisfying f\_U(A\textasciicircum{}par + g\_U(A\textasciicircum{}par)) = 0, where f\_U(A) = (1/2)*(((J + A\textasciicircum{}dagger) bold-dot U\textasciicircum{}dagger) bold-dot (U bold-dot (J + A)) - J), every tangent space T\_U calU is the image of L\_U(I + Dg\_U(0)): icalH -> calX, and omega\_U(Z) = (L\_U\textasciicircum{}\{-1\} Z)\textasciicircum{}par : T\_U calU -> icalH is a global C\textasciicircum{}1 bundle trivialization with distortion at most 1 + C\_ch*epsilon\_r, satisfying omega\_\{cU\}(cZ) = omega\_U(Z) and omega\_U(iU) = iJ for every U in calU, Z in T\_U calU, and c in U(1). \steptag{validated} \\[6pt]
\step{1.1} The external lem-stage1-maurer-cartan-trivialization supplies universal constants Cbar >= 1 and kappabar in (0,1/2] for which its stated uniform conclusion holds. \steptag{validated} \\[4pt]
\step{1.2} Choose the target constants C\_ch\textasciicircum{}0 := Cbar and kappa\_ch\textasciicircum{}0 := kappabar; these are universal and satisfy C\_ch\textasciicircum{}0 >= 1 and kappa\_ch\textasciicircum{}0 in (0,1/2]. \steptag{validated} \\[4pt]
\step{1.3} Fix arbitrary W, algebra, delta, and family g satisfying all hypotheses of node 1. The guard puts the algebra and delta within the scope of lem-stage1-maurer-cartan-trivialization. Without identifying g with the distinguished graph family mentioned there, the zero equation implies directly that the image of L\_U(I+Dg\_U(0)) is contained in T\_U calU; the external bundle trivialization gives equality of dimensions, hence equality of these spaces. The external omega, distortion, and equivariance conclusions then apply to this same tangent bundle, with preliminary distortion bound 1+Cbar*epsilon\_r. \steptag{validated} \\[6pt]
\step{1.3.1} From C\_ch*(epsilon\_r + delta) <= kappa\_ch, C\_ch >= Cbar, 0 < kappa\_ch <= kappabar, and epsilon\_r + delta >= 0, one has Cbar*(epsilon\_r + delta) <= C\_ch*(epsilon\_r + delta) <= kappa\_ch <= kappabar, which is precisely the smallness guard of lem-stage1-maurer-cartan-trivialization. \steptag{validated} \\[4pt]
\step{1.3.2} Conditional uniqueness bridge (and no more): suppose, in addition to the registered hypotheses, that the distinguished graph family gbar referred to in lem-stage1-maurer-cartan-trivialization is defined on B\textasciicircum{}\{icalH\}\_\{2delta\}(0), takes values in B\textasciicircum{}\{calH\}\_\{2delta\}(0), and satisfies f\_U(A\textasciicircum{}par + gbar\_U(A\textasciicircum{}par)) = 0 for every U and A\textasciicircum{}par. Then the target family's uniqueness hypothesis implies g\_U(A\textasciicircum{}par) = gbar\_U(A\textasciicircum{}par) for every U and A\textasciicircum{}par. The sole registered external does not supply this additional antecedent, so this node does not assert the unconditional equality g = gbar. \steptag{validated} \\[6pt]
\step{1.3.2.1} Fix arbitrary U in calU and A\textasciicircum{}par in B\textasciicircum{}\{icalH\}\_\{2delta\}(0), and assume the additional antecedent stated at node 1.3.2. Then gbar\_U(A\textasciicircum{}par) belongs to B\textasciicircum{}\{calH\}\_\{2delta\}(0) and solves f\_U(A\textasciicircum{}par + X) = 0. By the target hypothesis, g\_U(A\textasciicircum{}par) is the unique X in that ball solving this equation. Therefore gbar\_U(A\textasciicircum{}par) = g\_U(A\textasciicircum{}par). Since U and A\textasciicircum{}par were arbitrary, gbar = g pointwise. This proves precisely the conditional bridge and does not claim that the registered external establishes its antecedent. \steptag{validated} \\[4pt]
\step{1.3.3} Because g and gbar are C\textasciicircum{}1 and equal pointwise on B\textasciicircum{}\{icalH\}\_\{2delta\}(0), Dg\_U(0) = Dgbar\_U(0) for every U. Applying lem-stage1-maurer-cartan-trivialization under the guard in node 1.3.1, and substituting these derivative equalities, yields every tangent formula and the global C\textasciicircum{}1 omega formula for g, with distortion at most 1 + Cbar*epsilon\_r; the identities omega\_\{cU\}(cZ)=omega\_U(Z) and omega\_U(iU)=iJ are unchanged because omega is the same displayed map. \steptag{validated} \\[4pt]
\step{1.3.4} Repair of challenge ch-e67e0f72419358f7: no properties of, or equality with, the distinguished graph family are needed. The guard gives epsilon\_r < 1/2 because C\_ch >= 1, delta > 0, and C\_ch*(epsilon\_r+delta) <= kappa\_ch <= 1/2. If X satisfies X\textasciicircum{}dagger bold-dot X=J, then the C-star lower bound gives ||X||\textasciicircum{}2 <= 1/(1-epsilon\_r). For every Y, approximate associativity and exact unitality give ||L\_\{X\textasciicircum{}dagger\}L\_XY-Y|| <= epsilon\_r||X||\textasciicircum{}2||Y|| <= epsilon\_r/(1-epsilon\_r)||Y||, with coefficient strictly below 1. Thus L\_\{X\textasciicircum{}dagger\}L\_X is invertible by the Neumann lemma, so L\_X is injective; finite dimensionality makes L\_X surjective, hence X has a right inverse. Therefore every exact zero of the displayed f\_U equation lies in calU. \steptag{validated} \\[4pt]
\step{1.3.5} Fix U in calU and B in icalH. Exact unitality and U\textasciicircum{}dagger bold-dot U=J give f\_U(0)=0, so uniqueness forces g\_U(0)=0. For all sufficiently small real t, set X\_B(t)=U bold-dot (J+tB+g\_U(tB)). The defining zero equation is exactly X\_B(t)\textasciicircum{}dagger bold-dot X\_B(t)=J because the involution reverses the product, and the preceding step gives the required right inverse, so X\_B(t) lies in calU and X\_B(0)=U. Differentiating at zero gives dX\_B/dt at zero = L\_U(B+Dg\_U(0)B). Hence im L\_U(I+Dg\_U(0)) is contained in T\_U calU. \steptag{validated} \\[4pt]
\step{1.3.6} By node 1.3.1, lem-stage1-maurer-cartan-trivialization applies and supplies the fiberwise linear bundle isomorphism omega\_U:T\_U calU -> icalH, omega\_U(Z)=(L\_U\textasciicircum{}\{-1\}Z)\textasciicircum{}par. Consequently dim T\_U calU=dim icalH. For Phi\_g(B)=L\_U(B+Dg\_U(0)B), the inclusion from the preceding step makes omega\_U(Phi\_g(B)) defined, and the displayed formula gives omega\_U(Phi\_g(B))=(B+Dg\_U(0)B)\textasciicircum{}par=B because Dg\_U(0)B is Hermitian. Thus Phi\_g is injective, so its image has dimension dim icalH. The inclusion between finite-dimensional spaces of equal dimension is equality: T\_U calU=im L\_U(I+Dg\_U(0)). \steptag{validated} \\[4pt]
\step{1.3.7} The tangent bundle identified in the preceding step is the same intrinsic tangent bundle on which the external theorem already supplies omega\_U(Z)=(L\_U\textasciicircum{}\{-1\}Z)\textasciicircum{}par as a global C1 bundle trivialization, with distortion at most 1+Cbar*epsilon\_r and identities omega\_\{cU\}(cZ)=omega\_U(Z), omega\_U(iU)=iJ. These conclusions do not mention the distinguished graph family. Therefore they hold together with the newly established tangent formula for the arbitrary family g, completing node 1.3 without the missing antecedent flagged by the verifier. \steptag{validated} \\[4pt]
\step{1.4} Since C\_ch >= C\_ch\textasciicircum{}0 = Cbar and epsilon\_r >= 0, the preliminary distortion bound 1 + Cbar*epsilon\_r is at most 1 + C\_ch*epsilon\_r; hence the transferred conclusions are exactly all conclusions required by node 1, for every admissible datum. \steptag{validated} \\[4pt]
\end{proofbox}

\end{document}
